% =====================================================================
% 06-agentic-rl.tex
% =====================================================================
\section{From Single Responses to Agentic Trajectories}\label{sec:agentic}

The temporal scope of optimization distinguishes single-turn PBRFT from
\emph{agentic RL}, where an LLM is optimized over an entire multi-step
trajectory~\cite{zhang2026landscape}. Although the reward and
optimization dimensions were historically studied in the single-step
regime, the frontier has moved to trajectory-level training. This section
reviews that transition, the capabilities RL is used to instill, and the
task domains and infrastructure that support it.

\subsection{Why Trajectories Change the Problem}

Extending the horizon from $T{=}1$ to $T{>}1$ reintroduces the classical
difficulties of RL that PBRFT had sidestepped. Rewards become sparse and
delayed, raising the \emph{credit-assignment} problem of deciding which
of many actions caused a success or failure. The action space expands to
include tool calls and environment actions, demanding
\emph{exploration}. The policy must also maintain \emph{state}, that is,
memory of prior observations, across the episode. Process rewards
(Section~\ref{sec:reward}) and group-relative algorithms
(Section~\ref{sec:optimization}) are partial answers to these
difficulties, but they remain the defining challenges of the agentic
regime.

\subsection{RL-Optimizable Agentic Capabilities}

Following the capability view of Zhang \textit{et~al.}~\cite{zhang2026landscape},
RL is increasingly used to convert static, prompt-engineered modules into
learned, adaptive behaviors:

\begin{itemize}
\item \textit{Reasoning.} Long Chain-of-Thought (CoT), once elicited by
prompting~\cite{wei2022cot,yao2023tot,besta2024got}, is now \emph{trained}
with verifiable rewards, as in DeepSeek-R1~\cite{guo2025deepseekr1}.

\item \textit{Planning.} Plan--execute--reflect loops can be optimized
end-to-end rather than hand-designed, though LLM planning still degrades
sharply with combinatorial complexity~\cite{valmeekam2024planning}.

\item \textit{Tool use.} From self-supervised API
learning~\cite{schick2023toolformer,qin2023toolllm,patil2023gorilla},
the field is moving toward RL that learns \emph{when} and \emph{which}
tool to call based on final-answer quality.

\item \textit{Memory.} Episodic, semantic, and procedural memory are
emerging as RL-controllable policies for what to store and
retrieve~\cite{sun2025lifelong}.

\item \textit{Self-improvement.} Agents that critique and revise their
own trajectories close the loop between generation and reward.
\end{itemize}

\subsection{Task Domains}

Agentic RL has been instantiated across several domains, almost always
with verifiable rewards. In \emph{code}, SWE-agent couples an
LLM to an agent--computer interface and is evaluated on real GitHub
issues via SWE-bench~\cite{yang2024sweagent,jimenez2024swebench,wang2024codeact};
unit-test pass/fail provides a natural reward, and performance on
SWE-bench Verified rose from roughly 12.5\% to over 50\% within a year.
In \emph{math}, formal and informal verifiers drive
training~\cite{shao2024deepseekmath,lightman2023verify}. \emph{Search}
and \emph{GUI} agents optimize multi-turn retrieval and interface
manipulation. Embodied agents such as Voyager show open-ended skill
acquisition, albeit largely without RL in the loop~\cite{wang2023voyager}.

\subsection{Multi-Agent Systems}

A further extension optimizes \emph{multiple} interacting LLM agents.
Surveys catalogue debate, hierarchical orchestration, and social
simulation as the main collaboration
patterns~\cite{talebirad2024multiagent,han2025collaboration,du2023debate,hong2023metagpt,park2023generative}.
Applying RL to jointly train communication and delegation protocols among
agents remains largely unexplored and is a promising frontier. Practical
deployments of such agents, for instance intelligent agents for
computer-network configuration~\cite{reiznautt2026minicurso} and LLM
agents for communication and network
management~\cite{boateng2024llmnsm}, illustrate the demand for
trajectory-level training in realistic, tool-rich settings.

\subsection{Environments, Frameworks, and Benchmarks}

Trajectory-level training needs infrastructure: interactive environments
that expose tools and produce rewards, scalable RL frameworks that
overlap generation and learning, and benchmarks that measure multi-step
competence rather than single answers~\cite{zhang2026landscape}. The
field has rapidly accumulated open-source environments and training
frameworks, but standardized, contamination-resistant benchmarks for
\emph{adaptive, multi-turn} behavior remain scarce, a gap we return to
in Section~\ref{sec:challenges}.
